% Product Management - Week 14 Cornell Final Notes (Spring 2026, Mon/Tue)
% Compile: pdflatex product-management-final-cornell-master-notes-mon-tue-spring2026.tex
\documentclass[10pt,a4paper]{article}
\usepackage[utf8]{inputenc}
\usepackage[T1]{fontenc}
\usepackage[margin=1.2cm]{geometry}
\usepackage{xcolor}
\usepackage{array}
\usepackage{longtable}
\usepackage{enumitem}
\setlist[itemize]{leftmargin=1.2em,itemsep=1pt,topsep=2pt}
\setlist[enumerate]{leftmargin=1.4em,itemsep=1pt,topsep=2pt}
\definecolor{titleblue}{RGB}{25,55,109}
\definecolor{softgray}{RGB}{245,247,251}
\setlength{\parindent}{0pt}
\setlength{\parskip}{0.35em}

\newcommand{\cornell}[3]{
\subsection*{#1}
\rowcolors{1}{softgray}{white}
\begin{longtable}{|p{0.22\linewidth}|p{0.74\linewidth}|}
\hline
\textbf{Cue / Questions} & \textbf{Notes} \\
\hline
#2 & #3 \\
\hline
\end{longtable}
}

\begin{document}
\begin{center}
{\Large\bfseries\color{titleblue} Product Management - Week 14 Final Cornell Master Notes}\\
{\small Spring 2026 (typical class days: Mon/Tue) | Consolidated from class flow, weekly notes, and midterm prep materials}
\end{center}

\section*{Performance and Focus Rules}
\begin{itemize}
  \item \textbf{Product target level:} A (8.5/10) = notable/bueno and close to sobresaliente/excelente.
  \item \textbf{Risk thresholds to monitor globally:} WIP failing band = 3.49--0; UCM failing band = 4.99--0; suspenso = 4.
  \item \textbf{Focus principle:} keep Product at A while reallocating extra study time to weaker courses.
\end{itemize}

\section*{Cornell Notes by Core Theme}

\cornell{1) Product Foundations and Value Logic}
{What problem is solved?\\
Who is the user?\\
Why this solution now?\\
What value is unique?}
{
\textbf{Core concept:} Product work starts with a clear value hypothesis, not with features.\\
\textbf{Business model lens:} connect value creation (user utility) and value capture (revenue/cost logic).\\
\textbf{Decision checkpoint:} avoid building when user pain and willingness-to-pay are weakly evidenced.\\
\textbf{Exam language:} ``The proposal is strong if it demonstrates problem-solution fit and a credible path to repeated value delivery.''\\
\textbf{Common mistakes:} confusing idea quality with market demand; over-indexing on UI before validating core need.
}

\cornell{2) User Research, Interviews, and Validation}
{What assumptions are tested?\\
What is evidence vs opinion?\\
How do interviews change priorities?}
{
\textbf{Interview purpose:} reveal real context, pains, alternatives, and decision criteria.\\
\textbf{Validation cycle:} assumption -> test -> evidence -> iteration.\\
\textbf{Good evidence:} observed behavior, conversion/retention movement, consistent user patterns.\\
\textbf{Weak evidence:} ``I like it'' feedback without commitment signals.\\
\textbf{Operational rule:} one validated insight must translate into one concrete product decision.
}

\cornell{3) Prioritization, Roadmaps, and Delivery Discipline}
{What gets built now?\\
What is deferred and why?\\
How does roadmap show strategy?}
{
\textbf{Prioritization logic:} maximize expected impact under resource constraints.\\
\textbf{Roadmap role:} communication of strategic sequence, not feature wishlist.\\
\textbf{Execution link:} backlog quality determines sprint quality; vague items create waste.\\
\textbf{Trade-off examples:} retention improvement may dominate new acquisition features if churn is structurally high.\\
\textbf{Exam phrasing:} ``Prioritization is justified by impact, effort, risk reduction, and strategic coherence.''
}

\cornell{4) MVP, Experimentation, and Product-Market Fit}
{What is minimum viable?\\
What metric proves learning?\\
When is fit credible?}
{
\textbf{MVP meaning:} smallest solution that yields reliable learning signal.\\
\textbf{Experiment design:} state hypothesis, success threshold, timeline, decision rule.\\
\textbf{PMF signs:} recurring usage, retention stability, pull from target segment.\\
\textbf{Anti-pattern:} scaling acquisition before retention baseline is validated.\\
\textbf{Course tie-in:} lean validation reduces sunk cost risk and accelerates iteration quality.
}

\cornell{5) GTM and Growth Thinking}
{Who is first target segment?\\
Which channel economics work?\\
How is growth sustained?}
{
\textbf{GTM frame:} segmentation -> positioning -> channel -> pricing -> onboarding.\\
\textbf{Growth quality:} sustainable growth requires both acquisition and retention economics.\\
\textbf{North-star discipline:} one primary value metric with supporting diagnostic metrics.\\
\textbf{Market adaptation:} message and channel should match segment behavior, not generic assumptions.\\
\textbf{Exam cue:} always connect GTM choices to measurable business impact.
}

\cornell{6) Value Chain and Value Creation Case (Spain-SF-Japan)}
{Who is the ``Kus'' figure in class?\\
How do primary vs support activities differ?\\
What is valuation vs value creation?\\
How did 15K become 15M in this case?\\
Why did scale require international structure?\\
What is the IP trade-off: patents vs secrecy?}
{
\textbf{``Kus'' figure:} likely Michael Porter, linked to value creation logic through the value chain framework.\\
\textbf{Porter's value chain:} Primary activities (creation, sale, transfer) plus support activities (infrastructure, HR, technology, finance) must align.\\
\textbf{Value creation vs valuation:} resources remain economically weak until organized into a usable offer; valuation is a price estimate, value creation is real utility delivered.\\
\textbf{Case setup:} brain-disease R\&D centered on a molecular-mass discovery (Spain/San Francisco context).\\
\textbf{Value jump:} early research-stage value around 15K rose to around 15M after proof, structure, and market path were established.\\
\textbf{Global scaling logic:} Spain had limited market/structure; a major Japanese corporation recognized larger demand (250M/year potential) and enabled commercialization capacity.\\
\textbf{IP strategy trade-off:} patents require technical disclosure; secrecy can preserve edge when imitation risk is high.\\
\textbf{Asset logic:} technology, licenses, and methods are necessary but insufficient; without channels and operations, even a high-potential asset stays under-monetized.\\
\textbf{Exam phrasing:} ``Scientific discovery has low economic value by itself; value emerges when the discovery is embedded in a full value chain and scaled to a real market.''
}

\cornell{7) Branding, Rationality, and Innovation Strategy}
{What is bounded rationality in managerial decisions?\\
How does asymmetric information change consumer behavior?\\
Why do brands become trust proxies in uncertain markets?\\
Why is patenting both protection and exposure?\\
How does open innovation create value capture opportunities?\\
What is the Xerox PARC lesson for PM and strategy?}
{
\textbf{Bounded rationality:} managers and consumers decide with incomplete information, limited time, and cognitive constraints.\\
\textbf{Rationality vs emotion:} analytical structure guides decisions, but emotion heavily influences final choices under ambiguity.\\
\textbf{Lemon-market logic:} when sellers know more than buyers (asymmetric information), buyers rely on brand to avoid low-quality offers.\\
\textbf{Branding effect:} strong brands reduce perceived risk; when technical information is low, emotional trust can dominate choice.\\
\textbf{Patent strategy trade-off:} patents protect IP but require technical disclosure, which can increase imitation or litigation exposure.\\
\textbf{Cross-licensing logic:} firms often exchange IP rights to reduce legal conflict and speed ecosystem-level innovation.\\
\textbf{Open innovation (Chesbrough):} firms create more value by combining internal capabilities with external knowledge, partners, and platforms.\\
\textbf{Standard-setting strategy:} giving away a layer (e.g., software) can establish industry standards and shift monetization to services, hardware, or complements.\\
\textbf{Xerox PARC paradox:} Xerox invented breakthrough technologies (GUI, mouse, laser-printing ecosystem) but captured limited value due to business-model mismatch.\\
\textbf{Strategic takeaway:} invention quality is not enough; monetization requires fit among information strategy, brand trust, channel design, and business model.
}

\cornell{8) Testing as Investing and Market Entry Strategy}
{Why is testing treated as investing?\\
How do network effects change launch strategy?\\
What is value creation vs monetization in platform models?\\
Which customer-success metrics matter most?\\
How do credibility, lobbying, and endorsements influence adoption?\\
What does the bootstrapping-to-scale path teach?}
{
\textbf{Testing as investing:} avoid all-in spending; each sprint is a staged investment that buys evidence before larger capital commitments.\\
\textbf{Lean logic:} formulate hypotheses, run small experiments, and pivot/continue using observed market behavior rather than internal optimism.\\
\textbf{Network effects (N+1):} platform value increases with each additional participant; early adoption strategy must prioritize density and interaction loops.\\
\textbf{Platform entry implication:} ecosystem dominance (e.g., a strong OS/user base) can become the practical ``north star'' for developer and channel priorities.\\
\textbf{Value vs monetization:} products may create huge user value (Google-style utility) while monetizing indirectly via ads, data, or complementary services.\\
\textbf{Revenue discipline:} useful products still fail without a clear payer, pricing logic, and repeatable capture mechanism.\\
\textbf{Customer-success metrics:} retention, subscription continuity, and satisfaction/review quality are stronger health indicators than one-time usage spikes.\\
\textbf{Credibility and influence:} brand endorsements, stakeholder communication, and policy/institutional access can accelerate trust and market entry.\\
\textbf{Bootstrapping story:} early low-paid or free work validates feasibility; once proof exists, larger partners can scale the model and multiply valuation.\\
\textbf{Innovation funnel:} Phase 1 research/molecular idea -> Phase 2 MVP/testing -> Phase 3 scaled commercialization with major partners.
}

\cornell{9) Value Management in the Modern Workforce}
{What is the PM role across functions?\\
How do value creation and value capture differ in execution?\\
Why must strategy separate user from customer?\\
What data capabilities are now mandatory?\\
Why are communication and collaboration core skills?\\
How is AI used in real corporate workflows?}
{
\textbf{PM role as bridge:} product managers connect technology, finance, operations, and brand decisions into one coherent execution path.\\
\textbf{Value creation vs capture:} creation delivers user and ecosystem benefit; capture secures sustainable revenue/profit from that value.\\
\textbf{User vs customer logic:} the user may differ from the payer, so product design and monetization must satisfy both decision-makers.\\
\textbf{Data requirement:} modern roles expect fluency in external market signals and internal performance data for prioritization and positioning.\\
\textbf{Communication premium:} teams prioritize collaboration and translation skills (explaining analysis/AI outputs clearly), not only specialist technical depth.\\
\textbf{AI in practice:} tools are used to consolidate large information sets (e.g., hundreds of records) into decision-ready insights for managers.\\
\textbf{Company examples:} Meta emphasizes strategic PM impact at platform scale; Figma combines product and AI with design-led value; Santander links pricing/financial analysis with customer feedback loops; Inditex ties brand demand to operations/manufacturing evaluation.\\
\textbf{Workforce takeaway:} high-impact managers are cross-functional generalists who convert raw data into actionable insight and measurable customer outcomes.
}

\cornell{10) Case Studies in Operational Excellence}
{How did McDonald's redesign its kitchen model?\\
Why do Ryanair/Southwest optimize turnaround time?\\
How does Zara convert store feedback into production action?\\
Why do brand signals (e.g., green vs red) matter operationally?\\
How do AI/data loops enable speed-based advantage?\\
What does ``dated ground'' imply across industries?}
{
\textbf{Core principle:} speed and efficiency are direct value drivers when they reduce waste, idle time, and mismatch with demand.\\
\textbf{McDonald's pivot:} shifted from batch/pre-made output to ``Made For You'' flow, improving freshness/customization while requiring tighter kitchen coordination.\\
\textbf{Operational design lesson:} cross-functional alignment (ops, technology, training, finance) is necessary to sustain process redesign at scale.\\
\textbf{Aviation model (Ryanair/Southwest):} planes earn while flying, so fast gate turnaround increases asset utilization and supports low-fare economics.\\
\textbf{Throughput logic:} minimizing non-productive ground time lowers effective cost per unit delivered (meal, seat, garment).\\
\textbf{Zara/Inditex model:} near-real-time store feedback informs rapid replenishment and short-cycle production decisions.\\
\textbf{Scarcity mechanism:} limited availability and fast rotation create urgency, increasing conversion and reducing markdown risk.\\
\textbf{Brand perception shift:} visual repositioning (red to green cues) can signal quality/health/environment values as market expectations evolve.\\
\textbf{AI and data loops:} closed feedback systems connect customer behavior to headquarters decisions, accelerating demand sensing and response.\\
\textbf{``Dated ground'' meaning:} value is lost when assets sit idle (planes on ground, burgers in bins, shirts in warehouses); value is created by fast flow through the chain.
}

\cornell{11) Disruptive Innovation and Technology S-Curves}
{What is the technology S-curve?\\
Why did Kodak stay on the old curve?\\
How does Nokia vs. Apple illustrate platform vs product?\\
When does ``worse'' technology win on customer value?\\
How does Uber exemplify new-market disruption?\\
Why do incumbents lose to disruptors?}
{
\textbf{S-curve logic:} early performance is slow (R\&D), growth steepens, then a plateau as limits or economics cap improvement.\\
\textbf{Disruptive innovation:} entrants with fewer resources can win by serving overlooked segments or new value networks until quality meets mainstream needs.\\
\textbf{Kodak / digital camera:} Kodak helped invent digital imaging but protected film cash flow; rivals jumped the new curve while the incumbent delayed cannibalization.\\
\textbf{Nokia vs Apple:} Nokia optimized hardware durability and early ``small computer'' framing; Apple prioritized UX, ecosystem, and platform (App Store + services).\\
\textbf{Performance vs what customers value:} early digital photos traded quality for storage, speed, and convenience; ``better specs'' did not match the buying job.\\
\textbf{Incumbent lock-in:} legacy architectures (e.g., dominant silicon-era stack) persist when switching cost and risk outweigh marginal efficiency gains---``good enough'' endures.\\
\textbf{Quantum computing:} possible next S-curve---expensive and niche today, concentrated among large players, but a potential break from current chip paradigms.\\
\textbf{Uber case:} Uber Black targeted high-end/near-limo demand; UberX scaled accessibility; advantage included the information layer (location, price, predictability), not only the vehicle.\\
\textbf{Why giants fail:} fear of cannibalization; over-weighting current large customers who reject immature tech; incumbent complexity vs disruptor simplicity and ``good enough'' entry offers.
}

\cornell{12) Cost Structure and Strategic Design}
{How do price and cost differ strategically?\\
What are fixed vs variable costs, with examples?\\
How do economies of scale change unit economics?\\
What is Fordism (Model T, line, standardization)?\\
When is specialization cheaper than generalization?\\
What does ``land grab'' imply for fixed-cost bets?\\
How does the university case illustrate break-even logic?\\
How should managers match design to expected market volume?}
{
\textbf{Price vs cost:} price is what the customer pays; cost is what the firm incurs. A low-cost structure does not always mean the lowest market price---margin and positioning still matter.\\
\textbf{Fixed costs:} do not vary with output in the short run (rent, insurance, heavy equipment, plant).\\
\textbf{Variable costs:} scale with production (materials, much of direct labor, incremental energy).\\
\textbf{Economies of scale:} spreading fixed cost over more units lowers average cost per unit---central to mass production in large markets.\\
\textbf{Fordism / Model T:} standardization (``any color as long as it's black'') and assembly-line task division raised speed and allowed lower-skill, repeatable labor.\\
\textbf{Specialization vs generalization:} narrow roles are cheaper at high volume; broader, multi-skilled work supports flexibility and customization but often costs more per effective unit.\\
\textbf{Land-grab metaphor:} like claiming territory in the American West---once you commit to large fixed assets, you need scale and speed to justify the sunk structure.\\
\textbf{University as high-fixed-cost model:} campuses and infrastructure sit largely fixed; low enrollment vs capacity pressures profitability unless volume (students/revenue) covers the base.\\
\textbf{Break-even logic:} if fixed costs dominate and volume is insufficient, losses persist until revenue crosses the break-even point.\\
\textbf{Mass customization:} technology and process design can add variety without giving up all mass-production efficiency.\\
\textbf{Manager's design rule:} millions of units -> invest in capital-intensive scale and specialization; hundreds or niche demand -> keep fixed costs low and favor flexible, generalist capacity.
}

\cornell{13) Business Model Canvas and Zara (Inditex) Case}
{What is the BMC for in practice?\\
Who is the Zara segment in this analysis?\\
What is the value proposition for skinny jeans?\\
How do channels, relationships, and revenue fit Zara?\\
What is the operational timeline (weeks)?\\
What is the simplified Who / Why / How logic?\\
What caution applies when using AI for case data?}
{
\textbf{BMC role:} nine-block map of customers, offer, delivery, partners, activities, resources, costs, and revenue---used to align theory with a real operating model.\\
\textbf{Customer segments (student analysis):} roughly ages 17--27, mainly urban (e.g., NYC, LA, Miami), wanting premium-look fashion without luxury prices.\\
\textbf{Value proposition (skinny jeans line):} scarcity (buy now or gone), trend-led ``outfit solutions,'' mid-range price (e.g., \$30--\$50) vs Gap/Levi's but higher fashion relevance.\\
\textbf{Channels:} frequent store replenishment (e.g., twice-weekly in transcript) plus strong e-commerce.\\
\textbf{Customer relationships:} fast rotation drives repeat visits; weak sellers are cut quickly rather than deep-discounted.\\
\textbf{Revenue logic:} emphasis on full-price sales; scarcity reduces dependence on perpetual markdowns.\\
\textbf{Speed to market:} design-to-store in weeks (transcript: ~5-week cycle) vs months for traditional apparel---enables response to fast-moving trends.\\
\textbf{Simplified managerial logic:} Who (urban young adults) -> Why (trend, fit/scarcity) -> How (limited quantities, constant refresh).\\
\textbf{Zara ``secret'' in one line:} they monetize \textit{currentness}---small batches test demand; winners scale; losers exit with less inventory risk.\\
\textbf{Tooling caveat:} when using AI or simulations to populate a BMC, watch for plausible-but-wrong outputs (``hallucinations'') on pricing or facts; verify against sources.
}

\cornell{14) BMC Nine Blocks, Agile Supply Chain, and Zara Positioning}
{How do the nine BMC blocks work as a value-capture machine for Zara?\\
What defines agile vs traditional apparel retail?\\
Why does ``response beat prediction'' matter?\\
How should AI-generated case data be reality-checked?}
{
\textbf{Framing:} the nine blocks are not isolated boxes; together they describe how Inditex converts trend signals into repeatable value capture.\\
\textbf{1. Customer segments:} urban-oriented younger adults (e.g., 17--27) seeking ``premium'' styling without luxury price tags.\\
\textbf{2. Value propositions:} scarcity + \textit{currentness} (buy now or miss it); trend-led outfit solutions.\\
\textbf{3. Channels:} flagship/high-traffic urban stores plus tightly integrated digital (app/online) aligned with replenishment rhythm.\\
\textbf{4. Customer relationships:} largely self-service browsing/purchase, but emotionally driven by FOMO and rotation (repeat visits).\\
\textbf{5. Revenue streams:} predominantly full-price, higher-margin sales vs discount-dependent competitors.\\
\textbf{6. Key activities:} rapid design cycles, responsive production, JIT-style flow, continuous sell-through and data analysis.\\
\textbf{7. Key resources:} proprietary logistics and distribution capability (e.g., centralized hub / ``Cube'' narrative in class).\\
\textbf{8. Key partnerships:} external trend intelligence and flexible local manufacturing / supplier network.\\
\textbf{9. Cost structure:} higher logistics and coordination cost offset by lower markdown waste and faster inventory turns.\\
\textbf{Agile vs traditional:} many retailers use 6--9 month lead times and clear excess via deep discounts; Zara's ~5-week style cycle favors reactivity.\\
\textbf{Strategic contrast:} \textit{response over prediction}---small batches test the market; scale what works.\\
\textbf{Data verification lesson:} AI or simulations may oversimplify pricing tier (e.g., ``mid-range''); market positioning is often \textbf{premium fast fashion} (above mass basics, below designer), supporting impulse vs investment framing.\\
\textbf{Exam phrasing:} ``BMC blocks must align: the operating model (speed, logistics, partnerships) must match the value proposition (scarcity/currentness) and revenue model (full price).''
}

\cornell{15) Design Thinking and IDEO Shopping Cart Case}
{What is IDEO's cultural philosophy for idea quality?\\
How does empathy differ from asking ``experts'' only?\\
Why use multidisciplinary and paired observation?\\
How do define, ideate, prototype, and test link to user journey?\\
What does ``fail often'' mean as investment?\\
Why did the cart stay unchanged for decades?}
{
\textbf{IDEO philosophy:} best idea wins credibility, not hierarchy; flat culture encourages contributors at any level to surface insights.\\
\textbf{Phase 1 (empathy / observation):} study real behavior in context (e.g., grocery stores, ``tribes'' of shoppers and parents) rather than relying only on expert opinion.\\
\textbf{Observation method:} paired observers reduce individual bias; design-in-anthropology spirit---watch what people do, not only what they say.\\
\textbf{Multidisciplinary teams:} mix linguist, biologist, psychologist, MBA, etc., so non-obvious pain points surface that a single discipline might miss.\\
\textbf{Phase 2 (define / narrow):} compress research into focused problem frames (e.g., ~12 ``points of drama''), such as steering behavior or parking-lot cart abandonment.\\
\textbf{Phase 3 (ideate \& prototype):} ``Fail often to succeed sooner''---rapid, cheap prototypes are learning investments, not vanity polish; wall-based visuals/post-its beat slow text-only debate.\\
\textbf{Phase 4 (test the journey):} prototype the full \textbf{user journey}, not only the object; ``store of the future''--style physical simulation tests end-to-end flow.\\
\textbf{Five-stage map:} Empathize -> Define -> Ideate -> Prototype -> Test (non-linear in practice; iterate loops).\\
\textbf{Strategic insight:} stagnation persists when users accept ``that's how it is''; breakthrough starts with ``why must it stay this way?''---as with carts unchanged for ~50 years until reframed.\\
\textbf{Exam phrasing:} ``Design thinking ties empathy, diverse lenses, fast prototyping, and journey testing into one evidence loop for desirability and feasibility.''
}

\cornell{16) Customer Discovery, Lean Execution, and Yami (Workshop)}
{How does design thinking plus lean apply to a real student project?\\
What are ideation ground rules?\\
What is the golden rule of customer discovery interviews?\\
Which questions work vs which bias the answer?\\
How do you avoid designing for yourself?\\
Why time-box prototypes?\\
How does AI support vs replace interview data?}
{
\textbf{Project context (Yami):} social app aimed at reducing student isolation---workshop moves from IDEO theory to \textbf{customer discovery interviews}.\\
\textbf{Iterative reality:} ambition after ideation still returns to test with real people; winner ideas must earn evidence, not only team enthusiasm.\\
\textbf{Ideation ground rules:} defer judgment early; welcome wild ideas; build with ``Yes, and\ldots'' to compound rather than veto.\\
\textbf{Customer discovery purpose:} learn jobs-to-be-done, pain, and context \textbf{before} solution-selling; discover unsolved problems, not feature wishlists.\\
\textbf{Forbidden patterns:} avoid ``What features do you want?,'' ``Would you use an app that\ldots?,'' ``Do you like this idea?''---they produce politeness, false positives, and solution theater.\\
\textbf{Strong prompts (behavior + emotion):} ``Tell me about the last time you\ldots;'' ``What was most frustrating about\ldots;'' ``How did you feel when\ldots;'' ``Tell me about a normal weekday at uni;'' ``Describe a recent moment you felt connected;'' ``What apps do you use to talk to people most?;'' ``When did you want to join something but didn't---what stopped you?''\\
\textbf{Designer bias check:} separate your own emotions from observed user behavior; evidence beats internal conviction.\\
\textbf{Prototyping under constraint:} low-fidelity, clock-bound builds (e.g., by Day 5) force faster learning and limit over-engineering before validation.\\
\textbf{MVP link:} smallest artifact that tests a specific hypothesis about pain, frequency, or willingness to act---not a feature-complete app.\\
\textbf{FOMO note:} scarcity and social products often lean on fear-of-missing-out; validate whether that matches real student behavior vs assumption.\\
\textbf{Eclectic teams:} mixed disciplines still apply from IDEO logic for a full 360-degree view of the problem in workshop settings.\\
\textbf{AI + research:} tools can consolidate notes and check consistency, but \textbf{first-person interview stories} remain primary evidence for product decisions.\\
\textbf{Exam phrasing:} ``Discovery interviews seek past behavior and felt friction; they never outsource strategy to hypothetical feature votes.''
}

\cornell{17) Product Design, Modular Engineering, and User Research}
{What does user-centric research require next week?\\
What is modular design and why use a shared platform?\\
How does reverse engineering support design decisions?\\
How does form relate to function (sugar example)?\\
What does the car-mirror case illustrate about value engineering?\\
Why must you understand ``architecture'' before breakthrough ideation?}
{
\textbf{Research mandate:} run real interviews beyond friends; observe expert or segment ``masters''; Zoom or in-person notes; prioritize \textbf{unmet needs} over stated wishlists.\\
\textbf{Modular design:} build on a \textbf{platform}---shared internal architecture (chassis, motherboard logic) so multiple products reuse core subsystems.\\
\textbf{Benefits:} faster variants, lower duplicate engineering, path to \textbf{mass customization} (perceived variety at mass-production economics).\\
\textbf{Reverse engineering:} dissect benchmarks (mirrors, dispensers, etc.) to see mechanism, manufacturing cost drivers, and improvement levers.\\
\textbf{Form vs function:} design must solve behavior and production constraints, not only aesthetics---efficiency and cost are part of the design problem.\\
\textbf{Sugar dispenser arc:} (1) traditional spoons/large sweets habits, costly at scale; (2) 1990s-style measured-pour dispensers---cleaner, dose control; (3) modern materials (plastic vs metal/glass) driven by feel, durability, and environmental/cost trade-offs.\\
\textbf{Automotive mirror case:} larger/red-style mirror vs compact black mirror---smaller aero-friendlier module on a \textbf{shared platform}, less material, easier manufacture while preserving safety function (value engineering narrative).\\
\textbf{Modular creativity (music analogy):} composing by recombining blocks/modules illustrates the same combinatorial logic in product architecture.\\
\textbf{Ideation bridge:} next step converts interview ``pains'' and benefits into solution concepts (``shooting for the star'') with evidence anchors.\\
\textbf{Strategic insight:} breakthrough ideas require grasping physical, economic, and psychological \textbf{constraints} of the user---the ``rules of the game''---before sketching the hero concept.\\
\textbf{Exam phrasing:} ``Platform modularity trades integration depth for speed and customization; value engineering preserves function while attacking cost and complexity.''
}

\cornell{18) Manufacturing Strategy: Job Shop vs Mass Production}
{How do small batch and mass batch differ on cost and variety?\\
General-purpose vs special-purpose equipment?\\
Labor skill vs assembly-line specialization?\\
What are the three inventory buckets?\\
Make-to-order vs make-to-stock trade-offs?\\
How does the guitar case illustrate design for manufacture?\\
How should process choice follow market demand?}
{
\textbf{Batch scale:} small batches raise unit cost but enable customization (job shop); large batches spread fixed cost and unlock \textbf{economies of scale} for price competition.\\
\textbf{Equipment:} general-purpose machines are flexible, lower capex, slower; special-purpose machines target one operation, high upfront cost, very fast at volume.\\
\textbf{Labor:} high-skill generalists command premium wages and adapt; specialized line workers repeat one task, cheaper, less flexible.\\
\textbf{Product-process fit (matrix idea):} process type should align volume/variety---no universal ``best'' line; mismatch wastes capital or forfeits margin.\\
\textbf{Scheduling:} agile chains (e.g., fast fashion) need rapid rescheduling; mass lines need \textbf{standardization} and steady cadence (takt/cycle rhythm).\\
\textbf{Inventory types:} raw materials; \textbf{work-in-process (WIP)}; finished goods---each ties to cash, space, and obsolescence (\textbf{carrying cost}).\\
\textbf{Make-to-order:} start production after demand is known---flexible, minimal finished-goods stock, longer customer wait.\\
\textbf{Make-to-stock:} build to forecast---fast fulfillment, risk of unsold finished inventory and markdowns.\\
\textbf{Guitar comparison:} hand-built path emphasizes craft, detail, general tools, slow throughput, high perceived value; mass path emphasizes \textbf{standardization}, dedicated equipment, identical units at target cycle time.\\
\textbf{Manager's rule:} match the factory to demand---dozens of unique SKUs vs tens of thousands of one SKU implies opposite equipment, labor, and inventory design.\\
\textbf{Exam phrasing:} ``Operations strategy is choosing the batch, equipment, labor, and inventory posture that fits the market's volume, variety, and speed requirements.''
}

\cornell{19) User Research: Student Social Connectivity (Yami Context)}
{What pain points appear in mass university communication?\\
How does friendship formation change after week one?\\
Where do bonds actually form?\\
What is the ``social script'' problem?\\
How do students discover events digitally vs experience them physically?\\
What design opportunity does this imply for Yami?}
{
\textbf{Research setting:} empathy/discovery phase---focus-group style insights on student connection (links to Yami project).\\
\textbf{Notification fatigue:} high volume of mass emails (e.g., multiple per day) trains students to ignore inboxes or skim only ``Urgent'' tags---important invites become \textbf{out of sight, out of mind}.\\
\textbf{Freshman-week window:} early days offer maximum openness; later, people cluster into \textbf{closed circles} (2--3 friends), raising barrier for newcomers or circle expansion.\\
\textbf{Rituals and events:} bonds form through shared experiences---clubs, fundraisers, study tables, large gatherings; major-specific clubs blend friendship with resume signaling.\\
\textbf{Social script fatigue:} repetitive small talk (name, major, hometown) creates early energy then \textbf{social drain}; superficial loops feel tiring fast.\\
\textbf{Digital vs physical:} discovery often via aggregated social accounts (e.g., club-event Instagram); meaningful connection still depends on \textbf{in-person} co-presence and shared activity.\\
\textbf{``Object permanence'' (social colloquial):} without daily visibility, peers drop off the mental map---relevant to re-engagement and lightweight touchpoints.\\
\textbf{Social energy arc:} (1) excitement and openness; (2) annoyance/fatigue with repetition; (3) retreat into small stable groups (\textbf{psychographic} rhythm).\\
\textbf{Core tension:} desire for new relationships collides with awkwardness outside structured routes (clubs, major events)---\textbf{friction} and perceived safety matter.\\
\textbf{Product implication:} reduce noise vs mass email; lower friction from \textbf{online discovery} to \textbf{shared experience}; break the tired script with better context, timing, or structured icebreakers (hypothesis to validate).\\
\textbf{Exam phrasing:} ``User discovery translates pains (fatigue, scripts, closed circles) into behavioral requirements before feature lists.''
}

\cornell{20) Ideation, Prototyping, and Yami Feature Concepts}
{What are the workshop rules during ideation?\\
Which personas emerged from research?\\
What top pains did the group prioritize?\\
Which prototype-level features address authenticity?\\
How can LLMs support the next step?\\
Why kill ``static profile'' small talk?}
{
\textbf{Workshop phase:} moves from synthesized research (personas/pains) into \textbf{ideation} and \textbf{low-fidelity prototypes} for Yami.\\
\textbf{Ideation rules:} no criticism during the ``pour out'' phase; \textbf{Yes, and\ldots} to build on peers; prioritize \textbf{volume} of ideas before filtering.\\
\textbf{Personas (examples):} \textit{Social connector}---naturally introduces others, needs leverage to scale introductions; \textit{Niche explorer}---seeks values/interests over surface traits; \textit{Habitual participant}---bonds through repeated structured contexts (classes, clubs).\\
\textbf{Top pains:} (1) \textbf{social masking}---effort to perform a false-safe persona; (2) \textbf{small-talk fatigue}---repetitive name/major scripts; (3) \textbf{structured demands}---pressure in formal/professional social settings.\\
\textbf{Feature concepts:} \textbf{anti-static profiles} refreshed every 1--2 weeks so identity reflects current mood/interests; \textbf{unfiltered Q\&A}---users answer their own deep prompts first to seed trust; \textbf{bucket-list matching} on future goals (e.g., travel); \textbf{social buffering}---bring a +1/friend into events to lower entry anxiety (ties to gamified ice-breakers as optional engagement layer).\\
\textbf{Customer journey lens:} map Discovery $\rightarrow$ traction $\rightarrow$ deeper conversation $\rightarrow$ commitment; features should reduce superficiality and increase shared-experience authenticity.\\
\textbf{Matching logic:} algorithmic suggestions from values, lists, and dynamic signals vs random swipe mechanics (hypothesis---validate with users).\\
\textbf{AI handoff:} LLMs can help translate sketches/sticky notes into draft \textbf{functional requirements} or app logic---human judgment still validates desirability and ethics.\\
\textbf{Strategic insight:} stale profiles force recycled conversations; forced refresh pushes talk toward \textbf{present reality}, attacking the ``small talk killer'' problem.\\
\textbf{Exam phrasing:} ``Ideation defers judgment; prototypes test desirability cheaply; features must trace back to named pains and personas.''
}

\section*{Finals Tile Insights (From Uploaded Images)}
\begin{itemize}
  \item The Eco-Stride responses show recurring feedback around \textbf{precision}: distinguish differentiation logic from cost logic with explicit evidence.
  \item QFD and design-thinking answers need cleaner \textbf{step-by-step causal logic} (customer pain -> technical requirement -> value impact).
  \item Value capture answers should prioritize \textbf{brand + patents + channel/distribution + scalable operations} as complementary assets.
  \item Scoring comments indicate structure quality matters: concise claims, direct justification, and fewer broad statements.
\end{itemize}

\section*{Summary Block (Cornell Bottom)}
\begin{itemize}
  \item Product success in this course is driven by evidence quality, prioritization clarity, and execution consistency.
  \item For finals, answer with structure: problem -> evidence -> decision -> metric -> expected outcome.
  \item Maintain A-level by writing analytical answers with explicit trade-offs and decision criteria.
\end{itemize}

\end{document}
